% ============================================================
%  Math Notes Template
% ============================================================
\documentclass[11pt, a4paper]{article}

% --- Mathematics (must load before newtxmath) ---
\usepackage{amsmath, amssymb, amsthm, mathtools}

% --- Fonts & Encoding (XeLaTeX) ---
\usepackage{fontspec}
\usepackage{libertine}                    % Linux Libertine text + Biolinum sans
\usepackage[libertine]{newtxmath}         % Matching math font (loads after amsthm)
\usepackage{inconsolata}                  % Monospace font
\usepackage{microtype}                    % Improved typography

% --- Page Layout ---
\usepackage[
  margin=2.5cm,
  headheight=14pt
]{geometry}

% --- Line Spacing ---
\usepackage{setspace}
\onehalfspacing

% --- Paragraph Spacing ---
\usepackage{parskip}

% --- Colours ---
\usepackage[dvipsnames]{xcolor}
\definecolor{defblue}{HTML}{1a5276}
\definecolor{thmgreen}{HTML}{196f3d}
\definecolor{exred}{HTML}{922b21}
\definecolor{linkblue}{HTML}{2471a3}

% --- TikZ ---
\usepackage{tikz}
\usetikzlibrary{
  arrows.meta,
  calc,
  decorations.markings,
  patterns,
  positioning,
  shapes.geometric,
  cd                                      % commutative diagrams
}

% --- Theorem Environments ---
\usepackage{thmtools}

\declaretheoremstyle[
  headfont=\bfseries\color{thmgreen},
  bodyfont=\normalfont,
  spaceabove=10pt,
  spacebelow=10pt,
]{thmstyle}

\declaretheoremstyle[
  headfont=\bfseries\color{defblue},
  bodyfont=\normalfont,
  spaceabove=10pt,
  spacebelow=10pt,
]{defstyle}

\declaretheoremstyle[
  headfont=\bfseries\color{exred},
  bodyfont=\normalfont,
  spaceabove=10pt,
  spacebelow=10pt,
]{exstyle}

\declaretheorem[style=thmstyle, numberwithin=section, name=Theorem]{theorem}
\declaretheorem[style=thmstyle, sibling=theorem, name=Lemma]{lemma}
\declaretheorem[style=thmstyle, sibling=theorem, name=Proposition]{proposition}
\declaretheorem[style=thmstyle, sibling=theorem, name=Corollary]{corollary}
\declaretheorem[style=defstyle, sibling=theorem, name=Definition]{definition}
\declaretheorem[style=exstyle,  sibling=theorem, name=Example]{example}
\declaretheorem[style=exstyle,  numbered=no,      name=Remark]{remark}

% --- Hyperlinks & Cross-references ---
\usepackage[
  colorlinks=true,
  linkcolor=linkblue,
  citecolor=linkblue,
  urlcolor=linkblue
]{hyperref}
\usepackage{cleveref}

% --- Headers & Footers ---
\usepackage{fancyhdr}
\pagestyle{fancy}
\fancyhf{}
\fancyhead[L]{\textsc{\nouppercase{\leftmark}}}
\fancyhead[R]{\thepage}
\renewcommand{\headrulewidth}{0.4pt}

% --- Enumeration ---
\usepackage{enumitem}

% --- Bibliography (optional, uncomment if needed) ---
% \usepackage[style=alphabetic, backend=biber]{biblatex}
% \addbibresource{refs.bib}

% ============================================================
%  Fill these in for each set of notes
% ============================================================
\newcommand{\notestitle}{Basic Algebraic Topology}
\newcommand{\notesauthor}{B. Arm}
\newcommand{\notesdate}{\today}

\title{\notestitle}
\author{\notesauthor}
\date{\notesdate}

% ============================================================
\begin{document}

\maketitle
\tableofcontents
\newpage

% ============================================================
%  Chapter 1: Introduction
% ============================================================
\section{Introduction}

\subsection{The Basic Problem}

\subsection{Fundamental Group}

\subsection{Function Spaces and Quotient Spaces}

\subsection{Relative Homotopy}

\subsection{Some Typical Constructions}

\subsection{Cofibrations}

\subsection{Fibrations}

\subsection{Categories and Functors}

% ============================================================
%  Chapter 2: Cell Complexes and Simplicial Complexes
% ============================================================
\section{Cell Complexes and Simplicial Complexes}

\subsection{Basics of Convex Polytopes}

\subsection{Cell Complexes}

\subsection{Product of Cell Complexes}

\subsection{Homotopical Aspects}

\subsection{Cellular Maps}

\subsection{Abstract Simplicial Complexes}

\subsection{Geometric Realization of Simplicial Complexes}

\subsection{Barycentric Subdivision}

\subsection{Simplicial Approximation}

\subsection{Links and Stars}

% ============================================================
%  Chapter 3: Covering Spaces and Fundamental Group
% ============================================================
\section{Covering Spaces and Fundamental Group}

\subsection{Basic Definitions}

\subsection{Lifting Properties}

\subsection{Relation with the Fundamental Group}

\subsection{Classification of Covering Projections}

\subsection{Group Action}

\subsection{Pushouts and Free Products}

\subsection{Seifert--van Kampen Theorem}

\subsection{Applications}

% ============================================================
%  Chapter 4: Homology Groups
% ============================================================
\section{Homology Groups}

\subsection{Basic Homological Algebra}

\subsection{Singular Homology Groups}

\subsection{Construction of Some Other Homology Groups}

\subsection{Some Applications of Homology}

\subsection{Relation between $\pi_1$ and $H_1$}

\subsection{All Postponed Proofs}

% ============================================================
%  Chapter 5: Topology of Manifolds
% ============================================================
\section{Topology of Manifolds}

\subsection{Set Topological Aspects}

\subsection{Triangulation of Manifolds}

\subsection{Classification of Surfaces}

\subsection{Basics of Vector Bundles}

% ============================================================
%  Chapter 6: Universal Coefficient Theorem for Homology
% ============================================================
\section{Universal Coefficient Theorem for Homology}

\subsection{Method of Acyclic Models}

\subsection{Homology with Coefficients: The Tor Functor}

\subsection{K\"unneth Formula}

% ============================================================
%  Chapter 7: Cohomology
% ============================================================
\section{Cohomology}

\subsection{Cochain Complexes}

\subsection{Universal Coefficient Theorem for Cohomology}

\subsection{Products in Cohomology}

\subsection{Some Computations}

\subsection{Cohomology Operations; Steenrod Squares}

% ============================================================
%  Chapter 8: Homology of Manifolds
% ============================================================
\section{Homology of Manifolds}

\subsection{Orientability}

\subsection{Duality Theorems}

\subsection{Some Applications}

\subsection{de Rham Cohomology}

% ============================================================
%  Chapter 9: Cohomology of Sheaves
% ============================================================
\section{Cohomology of Sheaves}

\subsection{Sheaves}

\subsection{Injective Sheaves and Resolutions}

\subsection{Cohomology of Sheaves}

\subsection{\v{C}ech Cohomology}

% ============================================================
%  Chapter 10: Homotopy Theory
% ============================================================
\section{Homotopy Theory}

\subsection{$H$-Spaces and $H_0$-Spaces}

\subsection{Higher Homotopy Groups}

\subsection{Change of Base Point}

\subsection{The Hurewicz Isomorphism}

\subsection{Obstruction Theory}

\subsection{Homotopy Extension and Classification}

\subsection{Eilenberg--Mac Lane Spaces}

\subsection{Moore--Postnikov Decomposition}

\subsection{Computation with Lie Groups and Their Quotients}

\subsection{Homology with Local Coefficients}

% ============================================================
%  Chapter 11: Homology of Fibre Spaces
% ============================================================
\section{Homology of Fibre Spaces}

\subsection{Generalities about Fibrations}

\subsection{Thom Isomorphism Theorem}

\subsection{Fibrations over Suspensions}

\subsection{Cohomology of Classical Groups}

% ============================================================
%  Chapter 12: Characteristic Classes
% ============================================================
\section{Characteristic Classes}

\subsection{Orientation and Euler Class}

\subsection{Construction of Stiefel--Whitney Classes and Chern Classes}

\subsection{Fundamental Properties}

\subsection{Splitting Principle and Uniqueness}

\subsection{Complex Bundles and Pontrjagin Classes}

% ============================================================
%  Chapter 13: Spectral Sequences
% ============================================================
\section{Spectral Sequences}

\subsection{Warm-Up}

\subsection{Exact Couples}

\subsection{Algebra of Spectral Sequences}

\subsection{Leray--Serre Spectral Sequence}

\subsection{Some Immediate Applications}

\subsection{Transgression}

\subsection{Cohomology Spectral Sequences}

\subsection{Serre Classes}

\subsection{Homotopy Groups of Spheres}

% \printbibliography

\end{document}
